\selectlanguage{ngerman}%
\frenchspacing%
\begin{ledgroupsized}[r]{120mm}%
\footnotesize%
\pstart%
\noindent%
\textbf{Überlieferung:}%
\pend%
\end{ledgroupsized}%
\begin{ledgroupsized}[r]{114mm}%
\footnotesize%
\pstart%
\parindent -6mm%
\makebox[6mm][l]{\textit{L}}%
Aufzeichnung:
LH XXXVII\;4 Bl.\;107.
Ein Blatt 8\textsuperscript{o},
Ränder beschnitten:
Papiererhaltungsmaßnahmen;
geringfügiger Textverlust durch Tintenfleck.
Eine Seite auf Bl.\;\href{\Scan{LH_37_04_107r}}{107\,r\textsuperscript{o}}.
Bl.\;\href{\Scan{LH_37_04_107v}}{107\,v\textsuperscript{o}} ist unbeschrieben.
\pend%
\end{ledgroupsized}%
\vspace{5mm}%
\begin{ledgroup}%
\footnotesize%
\pstart%
\noindent%
\textbf{Datierungsgründe:}
\pagelabel{60261_dat}
Im Mittelpunkt der undatierten notizartigen Aufzeichnung N.\;\ref{60261} \textit{Funiculus Lini} steht die Erklärung, die der Jesuit Francis Line (Franciscus Linus)\indexN{\textls{Line} (Linus), Francis 1595\textendash1675} in seiner Abhandlung \textit{De corporum inseparabilitate} (London 1661)\cite{00072} über das Experiment abgibt, mit dem Gilles Personne de Roberval\indexN{\textls{Roberval} (Robervallius), Gilles Personne de 1602\textendash1675} den Zusammenhang zwischen dem Gewicht der Luft und ihrem elastischen Verhalten nachgewiesen hat:
Eine leere (Karpfen-)Blase dehne sich in einem geschlossenen Raum, der durch absinkenden Quecksilber luftentleert werde, wieder aus (vgl. %
ebd., cap.\;XII: S.\;55--58).\cite{00072}
Dieses Phänomen, das später oftmals in verbesserten und variierten Experimenten beobachtet wurde, u.\,a. von Robert Boyle\indexN{\textls{Boyle} (Boylius, Boyl), Robert 1627\textendash1691} (\textit{New experiments physico- mechanicall}, Nr.\;4: Oxford 1660, S.\;45--50; \textit{BW}~I, S.\;173--175),\cite{00015}\cite{00017}
will Francis Line\indexN{\textls{Line} (Linus), Francis 1595\textendash1675} erklären, ohne der Luft ein elastisches Verhalten zuzuschreiben:
Die Blase schwelle im luftentleerten Raum nicht deshalb, weil einige in ihr gebliebene Luft sich ausdehne, sondern deshalb, weil ihre Oberfläche von verdünntem Quecksilber, der den luftentleerten Raum besetze und sich wie eine sich zusammenziehende "Schnur" verhalte, gespannt werde (vgl. \textit{De corporum inseparabilitate}, cap.\;XII: S.\;58--60).\cite{00072}
Hiermit glaubte Line, das Phänomen -- anders als Roberval\indexN{\textls{Roberval} (Robervallius), Gilles Personne de 1602\textendash1675} und seine Anhänger -- verständlich machen zu können, ohne die herkömmlich peripatetische Annahme der Unmöglichkeit eines leeren Raumes in Frage zu stellen.
\pend%
\pstart%
Dieser Erklärung entgegnet Leibniz im Hauptteil der Aufzeichnung N.\;\ref{60261} zwei Gedankenexperimente, die sie empirisch widerlegen sollten.
Es ist aber bemerkenswert, dass er in den ersten Textzeilen Lines\indexN{\textls{Line} (Linus), Francis 1595\textendash1675} Ansatz ungenau wiedergibt:
Leibniz zufolge besteht der für das Schwellen der Blase verantwortliche \textit{funiculus} nicht, wie Line beteuert, aus verdünntem Quecksilber, sondern aus der "umgebenden Luft", die sich zusammenzieht, während der Raum entleert wird.
Diese Ungenauigkeit könnte bedeuten, dass Leibnizens Referat nicht auf eine direkte Lektüre zurückgeht, sondern möglicherweise auf eine vermittelnde, inakkurate Quelle.
\pend%
\pstart%
Hinweise auf Lines\indexN{\textls{Line} (Linus), Francis 1595\textendash1675} \textit{De corporum inseparabilitate}\cite{00072} finden sich im Leibniz-Nachlass bereits in Schriften der Mainzer\indexO{Mainz} Jahre (vgl. etwa \textit{Hypothesis physica nova} §\;59, Mainz 1671; \textit{LSB} VI,\;2 N.\;40, S.\;252.26)\cite{00256} sowie der frühen Pariser\indexO{Paris} Jahre (vgl. etwa \textit{LSB} VIII,\;1 N.\;\;40, S.\;326.9--18;\cite{01265} N.\;41, S.\;340.16-22).\cite{01190}
Bei keiner dieser früheren Erwähmungen lässt sich aber mit Sicherheit behaupten, dass sie von einer direkten Lektüre herrühren.
Ähnlich verhält es sich auch mit der Aufzeichnung N.\;\ref{60261} \textit{Funiculus Lini}, die somit ebenfalls in diesen frühen Jahren entstanden sein könnte.
Das kleine Papierformat sowie der Schreibduktus deuten aber nicht auf eine Abfassung vor Leibnizens Hannoveraner\indexO{Hannover} Zeit hin, d.h. vor 1677.
Eine frühere Entstehung ist trotzdem durchaus möglich.
\pend%
\pstart%
Ob und wann Leibniz Lines\indexN{\textls{Line} (Linus), Francis 1595\textendash1675} Abhandlung \textit{De corporum inseparabilitate}\cite{00072} gelesen hat, ist nach heutigem Wissensstand unklar.
Die umfangreichere Aufzeichnung \textit{Pro pressione aeris contra funiculum} (N.\;\ref{60206} im vorliegenden Band) berichtet ausführlicher -- wenngleich nicht mit viel größerer Genauigkeit -- als die Aufzeichnung N.\;\ref{60261} über Lines\indexN{\textls{Line} (Linus), Francis 1595\textendash1675} alternative Auffassung des Verhaltens der Luft und insbesondere über dessen Erklärung von Torricellis\indexN{\textls{Torricelli} (Torricellius), Evangelista 1608\textendash1647} berühmtem Experiment.
Dieser Umstand lässt die Vermutung zu, dass die Aufzeichnung N.\;\ref{60261} möglicherweise vor der Aufzeichnung N.\;\ref{60206} enstanden ist und dass Leibniz in der Zwischenzeit Gelegenheit gehabt haben könnte, sich weiter (wohl eher aus zweiter Hand) mit Lines\indexN{\textls{Line} (Linus), Francis 1595\textendash1675} \textit{Funiculus}-Lehre zu beschäftigen.
Da die Aufzeichnung N.\;\ref{60206} auf die Zeitspanne 1683--1684 datierbar ist (siehe die ausführliche Begründung auf S.\;\pageref{60206_dat}), dürfte auch N.\;\ref{60261} spätestens im Laufe des Jahres 1684 abgefasst worden sein.
Eine spätere Entstehung ist nicht ausgeschlossen, aber eher unwahrscheinlich.
Hieraus ergibt sich die für die Aufzeichnung N.\;\ref{60261} \textit{Funiculus Lini} vorgeschlagene Gesamtdatierung: 1677\;(?) bis 1684\;(?).
\pend%
\end{ledgroup}%
\selectlanguage{latin}%
\frenchspacing%
\vspace{8mm}%
\normalsize%
\pstart%
\noindent%
\lbrack%
\href{\Scan{LH_37_04_107r}}%
{107\,r\textsuperscript{o}}%
\rbrack%
\hfill%
Funiculus Lini%
\indexS{funiculus Lini}
\hfill%
\phantom{\lbrack107\,r\textsuperscript{o}\rbrack}%
\pend%
\vspace{0.5em}%
\pstart%
\noindent%
\edtext{%
Ait Linus%
\indexN{\textls{Line} (Linus), Francis 1595\textendash1675}
\edtext{%
vesicam antea flaccidam%
\indexS{vesica flaccida}
in recipiente exhausto%
\indexS{recipiens exhaustus aere}
inflari,%
\indexS{vesica inflata}%
}{%
\lemma{vesicam}%
\Bfootnote{%
\textit{(1)}~in re %
\textit{(2)}~antea flaccidam in recipiente exhausto %
\textit{(a)}~distendi, %
\textit{(b)}~inflari,%
~\textit{L}%
}}
ob vim aeris%
\indexS{vis aeris circumstantis}%
\indexS{aer circumstans}
\edtext{%
circumstantis,
qui nimium extensus%
\indexS{aer extensus}
dum se contrahere nititur,%
\indexS{aer nitens contrahere se}
funiculo suo%
\indexS{funiculus Lini}%
}{%
\lemma{circumstantis,}%
\Bfootnote{%
\textit{(1)}~funiculo suo %
\textit{(2)}~qui nimium %
\lbrack...\rbrack\ contrahere nititur, %
\textit{(a)}~funicu %
\textit{(b)}~funiculo suo%
~\textit{L}%
}}
vesicam attrahit,%
\indexS{vesica flaccida}
ut ipsa contrahentis sese locum impleat.%
\indexS{locus aeris contrahentis}%
\indexS{aer contrahens se}%
}{%
\lemma{Ait \lbrack...\rbrack\ impleat}%
\Cfootnote{%
Siehe
\textsc{F.\;Line},
\textit{De corporum inseparabilitate},
cap.\;XII
(London 1661,
S.\;55--60, bes. S.\;58\,f.).%
\cite{00072}
Leibnizens Wiedergabe ist ungenau.%
}}%
\pend%
\pstart%
Sed
\edtext{%
huic explicationi%
\indexS{explicatio}
sequens experimentum%
}{%
\lemma{huic}%
\Bfootnote{%
\textit{(1)}~sententiae\indexS{sententia} %
\textit{(2)}~explicationi sequens %
\textit{(a)}~\textlangle\textemdash\textrangle\ %
\textit{(b)}~experimentum%
~\textit{L}%
}}
opponi%
\indexS{experimentum oppositum}
\edtext{%
potest.
\edlabel{LH_37_04_107_107r_instantiaprima-1}%
Sit vesica instructa epistomio,%
\indexS{vesica instructa epistomio}%
\indexS{epistomium in vesica}
quod extrorsum
ubi inflata est%
\indexS{vesica inflata}
circumagendo,
intus aperiatur.%
}{%
\lemma{potest.}%
\Bfootnote{%
\textit{(1)}~Vesica intra recipientem inflata artificio aliquo\indexS{recipiens} %
\textit{(a)}~dirumpatur %
\textit{(b)}~pungatur, ut fiat in ea exiguum foramen, %
\textit{(2)}~Sit vesica \lbrack...\rbrack\ intus aperiatur.%
~\textit{L}%
}}
Statim subsidebit,
nec clauso Epistomio denuo inflabitur,%
\indexS{epistomium in vesica}
quod ob praecedentem rationem fieri%
\indexS{ratio}
\edlabel{LH037_04_107_107r_saltus-1}%
deberet.%
\edlabel{LH_37_04_107_107r_instantiaprima-2}%
\edtext{}{%
\xxref{LH037_04_107_107r_saltus-1}{LH037_04_107_107r_saltus-2}%
\lemma{deberet.}%
\Bfootnote{%
\textit{(1)}~Ratio etiam explicationi Linianae in eo contraria est%
\indexS{explicatio Liniana} %
\textit{(2)}~Objici etiam Lino potest,%
~\textit{L}%
}}
\pend%
\pstart%
Objici etiam Lino potest,%
\indexN{\textls{Line} (Linus), Francis 1595\textendash1675}%
\edlabel{LH037_04_107_107r_saltus-2}
quod si
\edtext{%
aer nimium extensus%
\indexS{aer extensus}
attrahendi%
}{%
\lemma{aer}%
\Bfootnote{%
\textit{(1)}~attrahendi %
\textit{(2)}~nimium %
\textit{(a)}~disten %
\textit{(b)}~extensus attrahendi%
~\textit{L}%
}}
virtute praeditus est,%
\indexS{virtus attrahendi}%
\indexS{aer praeditus virtute attraendi}
utique aer intra vesicam%
\indexS{aer intra vesicam}%
\indexS{vesica inflata}
aeri
\edtext{%
externo%
\indexS{aer externus}
extendenti ipsum
ac se contrahenti%
\indexS{aer contrahens se}
resistet.%
}{%
\lemma{externo}%
\Bfootnote{%
\textit{(1)}~contrahenti resistet %
\textit{(2)}~extendenti ipsum %
\lbrack...\rbrack\ contrahenti resistet.%
~\textit{L}%
}}
Responderi tamen potest pro Lino%
\indexN{\textls{Line} (Linus), Francis 1595\textendash1675}%
\indexS{responsio pro Lino}
eo minus resistere
quo minus jam extensus est.%
\indexS{aer extensus}
Sed
\edtext{%
prior instantia mea%
\indexS{instantia}%
}{%
\lemma{prior instantia mea}%
\Cfootnote{%
\refpassage{LH_37_04_107_107r_instantiaprima-1}{LH_37_04_107_107r_instantiaprima-2}.%
}}
refutari credo non potest.
\pend%
\selectlanguage{latin}%
\frenchspacing%
\normalsize%
\newpage%
